% -*- mode : latex; coding: utf-8 -*-
\chapter{Introduction}

DRAM has long been a popular choice for main memory in a computing system, 
and its cells are becoming smaller and closer to each other due to technology scaling. 
This makes DRAM vulnerable to inter-cell interference. \textit{RowHammer} (RH) attack 
is one type of such interference, where a burst of ACT command to a single DRAM row causes
 a charge leak in adjacent rows. 

This problem is a major threat to DRAM reliability, and thus many solutions have been proposed 
by both industry and academia. Recent protection schemes are based on architectural solutions, 
which implement hardware logic to detect hammering rows and handle them proactively. 
This class of mitigation scheme can be divided into two parts: probabilistic \cite{FlippingBits} and 
counter-based \cite{Graphene}, \cite{Hydra}, \cite{TWICE}, \cite{CBT}. 
One of the most famous probabilistic schemes is PARA \cite{FlippingBits}, which refreshes adjacent rows 
with small predefined probability. Although it doesn't require complex hardware logic, 
it invokes a substantial amount of unnecessary refreshes to ensure the required security level. 
However, counter-based protection only refreshes victim rows deterministically, making it an 
attractive choice. \textit{Graphene} \cite{Graphene} is best known counter-based protection 
scheme which uses Misra-Gries streaming algorithm \cite{MISRA-GRIES} to track hammering rows. 
It accurately tracks a set of rows that were activated more than a certain threshold from incoming commands.

The metric used to characterize how a DRAM chip is vulnerable to RH attack is $Flip_{th}$, 
which indicates the minimum number of row activations (ACT) required to induce a bit-flip in adjacent rows. 
Recent studies show that $Flip_{th}$ has dropped from 139K (DDR3) to 50K in DDR4 and 4.8K for LPDDR4 
\cite{FlippingBits}, \cite{TRRespass}, \cite{Revisiting}. As technology continues to scale, 
$Flip_{th}$ will likely decrease further, even lower than 1K. \textit{Hydra} \cite{Hydra} is a 
counter-based mitigation scheme targeting such ultra-low threshold with acceptable storage overhead. 
It uses the SRAM counter to track ACTs of the DRAM row group and uses the DRAM counter to track 
each row when the group counter reaches a certain threshold. 

Although promising, all of these mitigation strategies are inefficient in various aspects because there 
are no public documents available for internal DRAM microarchitecture. In this paper, I point out some 
optimization ways of \textit{Graphene} and \textit{Hydra} if such internal DRAM topologies are available. 
I implemented \textit{Graphene}, \textit{Hydra}, and optimized version of each in Macsim \cite{MACSIM} and 
Ramulator \cite{RAMULATOR}. Using these simulators, I evaluated storage overhead, the number of additional refreshes, and 
performance in synthetic and normal workloads such as SPEC CPU 2017. Below are the key contributions of this work.
% TODO: poor listing!
\begin{itemize}
  \item Graphene optimization: As DRAM's internal row mapping is not known, 
  the hardware tracker doesn't know the adjacent rows of a hammered row. So tracking algorithm of 
  Graphene always has to assume the attack is done as \textit{double-sided}, where one victim row is adjacent to two hammered rows. 
  I tracked victim rows instead of hammering rows and analyze the benefits and security of the modified algorithm.
  \item \textit{Hydra} optimization: Recent comprehensive analysis of DDR4 device showed that due to design and manufacturing factors, 
  $Flip_{th}$ variation exists across rows in DRAM. In the optimized version of  Hydra, 
  I stored the encoded $Flip_{th}$ value in each row's counter in DRAM.
  \item Evaluation: I implemented the above optimization points in a cycle-level  
  simulator and evaluated storage overhead, additional refresh, and performance. 
  In Graphene, an additional refresh was reduced by 50\% in attack patterns and average performance overhead was reduced by 65.7\%
  in normal workloads with a negligible increase in table size. In the case of Hydra, an additional refresh was reduced by 91.4\% in attack patterns and performance overhead was reduced by 71.7\% in normal workloads.
\end{itemize}

In this paper, I first illustrate the background of DRAM structure and RH mitigation 
schemes of Graphene and Hydra in Chapter \ref{sec:Backgrounds}. In Chapter \ref{sec:Optimization}, 
I explain improvement points and implementation details with security analysis. After that, I compare storage overhead, 
the number of an additional refreshes, and the performance (IPC) of each scheme's baseline and optimized version 
in Chapter \ref{sec:Evaluation}. Finally, Chapter \ref{sec:related-work} summarizes related work in progress, and
Chapter \ref{sec:Conclusion} concludes this paper and suggests possible future works.
